\section{A finite recognition interface}
For a partition $\pi$ of a finite set $Q$, write
$\Y(\pi)=\prod_{B\in\pi}\Sym(B)$, as a subgroup of $\Sym(Q)$.
Let $(h_i)_{i\in I}$ be a finite family in $\Sym(Q)$ and let
$H_S=\langle h_i:i\in S\rangle$ for $S\subseteq I$.
Fix $d\in Q$. Suppose each $S$ is assigned a partition $\pi_S$ of $Q$.
If $x\notin Q$, let $\pi_S^x$ be the partition of $Q\cup\{x\}$
obtained by adjoining $x$ to the block containing $d$.

We impose the following finite conditions, for every $S,R\subseteq I$:
\begin{align}
H_S\cap H_R&=H_{S\cap R}, &
H_S\cap\Y(\pi_R)&=H_{S\cap R},\label{eq:finite1}\\
\pi_S\wedge\pi_R&=\pi_{S\cap R}, &
\langle H_S,(d\ x)\rangle&=\Y(\pi_S^x).
\label{eq:finite2}
\end{align}
Here the meet is common refinement. In particular, taking $S=R$ in
\eqref{eq:finite1} gives $H_S\leq\Y(\pi_S)$. These hypotheses concern
subgroups on $|Q|$ or $|Q|+1$ letters; they do not assume the desired
intersection property on any larger degree.

Let $T$ be a finite tree meeting $Q$ in exactly $\{d\}$, with at least
one edge, and suppose $d$ is a leaf of $T$. Put
$\Omega=Q\cup V(T)$. Alongside the $h_i$, take the transposition
$t_e$ exchanging the endpoints of each edge $e$ of $T$. All these
permutations fix the unspecified letters.

\begin{theorem}\label{thm:attachment}
Under \eqref{eq:finite1}--\eqref{eq:finite2}, the family
$(h_i)_{i\in I},(t_e)_{e\in E(T)}$ has the full subset intersection
property. If the $h_i$ are distinct nonidentity involutions, then so is
this enlarged family and it is a C-group marking of its generated group.
If moreover $\pi_I=\{Q\}$, that group is $\Sym(\Omega)$.
\end{theorem}

\begin{proof}
For $F\subseteq E(T)$, let $D_F$ be the vertex set of the component
containing $d$ in $(V(T),F)$; the other component sets will be denoted
$C$. The tree transpositions on a connected component generate its full
symmetric group: along the unique paths, conjugating endpoint
transpositions produces every transposition from one fixed vertex.
Thus
\begin{equation}\label{eq:all-subsets}
\langle h_i\ (i\in S),t_e\ (e\in F)\rangle
= H_S[D_F]\times\prod_{C\ne D_F}\Sym(C),
\end{equation}
where $H_S[\{d\}]=H_S$. If $D_F\ne\{d\}$, then $H_S[D_F]$
is the Young subgroup obtained from $\pi_S$ by adjoining
$D_F\setminus\{d\}$ to its $d$-block. Indeed, choose any
$x\in D_F\setminus\{d\}$. The subgroup $\Sym(D_F)$ contains
$(d\ x)$, so \eqref{eq:finite2} supplies every transposition within
each unchanged core block and within the enlarged $d$-block on
$Q\cup\{x\}$. Conjugation by $\Sym(D_F)$ supplies the same
transpositions for every point of $D_F$. Conversely all the original
generators preserve the claimed blocks. This proves \eqref{eq:all-subsets}.

We now intersect the groups for $(S,F)$ and $(R,E)$. Intersections of
connected subtrees of a tree are connected or empty. More precisely,
two vertices belong to one component in both forests if and only if
the unique path between them belongs to $F\cap E$. Hence common
refinement of their component partitions is the component partition
of $F\cap E$. This proves the classical forest-intersection argument used throughout
the present construction; it is also the argument of our earlier UB.F.

If $D_F=D_E=\{d\}$, the core intersection is $H_{S\cap R}$
by \eqref{eq:finite1}; all other factors follow from the forest
argument. Suppose exactly one of $D_F,D_E$ is a singleton, say
$D_F=\{d\}$. Every permutation in the first group preserves $Q$
setwise and acts there in $H_S$. Its additional membership in the
second group restricts this action exactly to $\Y(\pi_R)$.
Equation \eqref{eq:finite1} gives $H_{S\cap R}$ on $Q$.
Off $Q$, the two forest component partitions have the refinement
just described; their intersections with the enlarged root component
are connected subtrees too. Consequently the remaining factors are
exactly the symmetric groups on the nonroot components of $F\cap E$.

Finally suppose both root components are nonsingletons. The unique
edge incident with $d$ lies in both $F$ and $E$. Therefore
$D_{F\cap E}=D_F\cap D_E$ is also nonsingleton. Both groups in
\eqref{eq:all-subsets} are now Young subgroups of $\Sym(\Omega)$.
Their intersection is the Young subgroup of the common refinement.
On the core this is $\pi_S\wedge\pi_R=\pi_{S\cap R}$;
the root block acquires exactly $D_F\cap D_E\setminus\{d\}$,
and all other blocks are the nonroot components of $F\cap E$.
This is precisely \eqref{eq:all-subsets} for $(S\cap R,F\cap E)$.
All subset intersections have been proved.

No tree transposition is supported on $Q$, and different edges give
different transpositions. Thus the distinctness assertion follows.
For the whole edge set, $D_F=V(T)$ is nonsingleton; if $\pi_I$
has one block, \eqref{eq:all-subsets} is $\Sym(\Omega)$.
\end{proof}

